\chapter{Fazit}
\label{chap:fazit}

Ausgangspunkt dieser Arbeit war die Frage, ob ein gezielt für Reflexion statt für schnelle Antworten
gestaltetes KI-System Studierende bei Prozessentscheidungen im Software Engineering anders
unterstützt als ein generisches KI-System. Dazu wurde ein Reflexionsbot konzipiert und entwickelt,
der Studierende durch sokratisches Fragen und adaptive Antworttiefe statt direkter Lösungen begleitet
(Kapitel~\ref{cha:concept}), und im Rahmen eines Mini-Experiments in der Lehrveranstaltung
\glqq Grundlagen des Software Engineering 1\grqq{} einer Vergleichsgruppe gegenübergestellt, die mit
einem generischen KI-System arbeitete. Die entstandenen Bearbeitungen und Chatverläufe wurden mittels
strukturierender qualitativer Inhaltsanalyse nach dem Reflexionsstufenmodell von Hatton und Smith
kodiert und mit den Ergebnissen eines Prä-Post-Surveys zusammengeführt (Kapitel~\ref{cha:methodik}).

Das zentrale Ergebnis fällt vorsichtiger aus, als die Ausgangserwartung nahelegte: In der
Reflexionsbot-Gruppe ist der Anteil der Bearbeitungen ohne eindeutig studentische Autorenschaft
geringer als in der Vergleichsgruppe (50\,\% gegenüber 80\,\%). Dieser Rohabstand ist jedoch teilweise
durch die Aufgabenvorgabe für die Vergleichsgruppe bedingt und fällt bei den gesichert KI-generierten
Fällen deutlich kleiner aus (Abschnitt~\ref{sec:ergebnisse}). Die direkten Lernmaße stützen keinen
Systemvorteil. Am objektiven Wissenstest lässt der Deckeneffekt zwischen den Systemgruppen keine
Aussage zu, und in der gezeigten Reflexionstiefe der studentischen Bearbeitungen liegt die
Vergleichsgruppe deskriptiv sogar höher, bei einer Fallzahl von drei gegenüber acht auswertbaren
Fällen, die keine Richtungsaussage trägt. Ein günstigeres Bild ergibt sich vor allem beim
selbsteingeschätzten Verständnis für Begründung und Abwägung, das in der Reflexionsbot-Gruppe stabil
bleibt oder leicht steigt, während es in der Vergleichsgruppe sinkt. Reflexionstiefe hing in den
vorliegenden Daten damit weniger vom eingesetzten Werkzeug ab als von der Bereitschaft der
Studierenden, sich auf den Dialog einzulassen. Ein erheblicher Teil der Reflexionsbot-Gruppe wich den
Rückfragen dennoch aus oder versuchte, das System zu direkten Antworten zu bewegen
(Abschnitt~\ref{sec:beantwortung-forschungsfragen}).

Die in Kapitel~\ref{chap:einleitung} formulierte Sorge, dass generische KI-Systeme das ungeprüfte
Übernehmen von Lösungen begünstigen, ist mit den vorliegenden Daten vereinbar. Für die Annahme, dass
ein didaktisch gezielt gestaltetes System dem entgegenwirkt, gilt das ebenfalls, ein Beleg ist sie
jedoch nicht. Der deutlichste Unterschied entsteht in einer Kennzahl, die durch die Aufgabenstellung
für die Vergleichsgruppe mitbedingt ist und zusätzlich durch eine Asymmetrie in der Verfügbarkeit der
Chatverläufe beeinflusst wird; bei den gesichert KI-generierten Fällen schrumpft der Abstand auf ein
Maß, das bei dieser Fallzahl nicht mehr interpretierbar ist (Abschnitt~\ref{sec:limitationen}). Für
die Praxis folgt daraus vor allem, dass die Gestaltung eines KI-Systems die Notwendigkeit nicht
ersetzt, Studierende zur inhaltlichen Auseinandersetzung zu motivieren. Konkrete Empfehlungen hierzu
finden sich in Abschnitt~\ref{sec:handlungsempfehlungen}.

Diese Arbeit deutet insgesamt darauf hin, dass didaktisch gestaltete KI-Systeme das Potenzial haben,
oberflächliche Nutzung generativer KI in der Hochschullehre zu reduzieren und -- bei Studierenden, die
sich darauf einlassen -- tiefere Reflexion zu begünstigen. Auf dem direkten Reflexionsstufenmaß zeigt
sich dieses Potenzial in den vorliegenden Daten allerdings nicht. Ob es tatsächlich zum Tragen kommt,
entscheidet sich zudem nicht allein am technischen Design des Systems, sondern daran, ob und wie sich
Studierende auf das Angebot zur Reflexion einlassen.
